\documentclass[a4paper,12pt]{article}
\usepackage[T2A]{fontenc}
\usepackage[utf8]{inputenc}
\usepackage[russian]{babel}
\usepackage{amsmath,amssymb,mathtools}
\usepackage[sfdefault,type1]{noto}
\usepackage{icomma}
\usepackage[a4paper,margin=19mm,top=22mm,bottom=20mm]{geometry}
\usepackage{microtype}
\usepackage{tikz}
\usetikzlibrary{calc,arrows.meta,positioning,shapes.geometric}
\usepackage{tkz-euclide}
\usepackage{booktabs,tabularx}
\usepackage{enumitem}
\usepackage{fancyhdr}
\usepackage{setspace}
\usepackage{xcolor}
\usepackage{hyperref}
\hypersetup{hidelinks,pdftitle={Чек-лист. Закон всемирного тяготения и сила тяжести},pdfauthor={Лёвин Артём Александрович}}

\definecolor{accent}{HTML}{166D72}
\definecolor{darkaccent}{HTML}{17484C}
\definecolor{textgray}{HTML}{383D40}
\definecolor{softblue}{HTML}{E9F1F1}
\color{textgray}
\onehalfspacing
\setlength{\parindent}{0pt}
\setlength{\parskip}{5pt}
\setlist[itemize]{leftmargin=7mm,itemsep=2pt,topsep=2pt}
\setlist[enumerate]{leftmargin=8mm,itemsep=4pt,topsep=4pt}
\newcommand{\checkline}{\ensuremath{\square}\,}
\newcommand{\key}[1]{\textcolor{darkaccent}{\textbf{#1}}}
\newcommand{\sect}[1]{\vspace{4pt}{\Large\bfseries\color{darkaccent}#1}\par\vspace{3pt}\textcolor{accent!60}{\rule{\linewidth}{0.6pt}}\vspace{2pt}}
\newcommand{\subsect}[1]{\vspace{3pt}{\large\bfseries\color{darkaccent}#1}\par}
\pagestyle{fancy}
\fancyhf{}
\fancyhead[L]{\small\textcolor{darkaccent}{Физика \textbullet\ 24.10.2026}}
\fancyhead[R]{\small\textcolor{darkaccent}{Володя Хачатурян}}
\fancyfoot[C]{\small\thepage}
\renewcommand{\headrulewidth}{0.4pt}
\renewcommand{\footrulewidth}{0pt}
\setlength{\headheight}{15pt}

\tikzset{
  flowterminal/.style={ellipse,draw=accent,line width=0.9pt,fill=softblue,align=center,text width=78mm,minimum height=12mm,inner sep=3mm,font=\small},
  flowstep/.style={rounded corners=2mm,rectangle,draw=accent!80,line width=0.9pt,align=center,minimum height=15mm,text width=105mm,inner sep=3mm,font=\small},
  flowbranch/.style={rounded corners=2mm,rectangle,draw=accent!80,line width=0.9pt,align=center,minimum height=25mm,text width=63mm,inner sep=2mm,font=\small},
  flowdecision/.style={diamond,aspect=2.8,draw=accent,line width=1pt,fill=softblue,align=center,text width=39mm,inner sep=1mm,font=\small},
  flowarrow/.style={-{Latex[length=3mm]},line width=0.9pt,draw=darkaccent}
}

\begin{document}
\begin{titlepage}
\thispagestyle{empty}
\begin{center}
\vspace*{29mm}
{\large\color{accent}\bfseries ФИЗИКА \textbullet\ 24 ОКТЯБРЯ 2026 ГОДА}\par
\vspace{22mm}
{\fontsize{30}{36}\selectfont\color{darkaccent}\bfseries Чек-лист}\par
\vspace{11mm}
{\fontsize{22}{29}\selectfont Закон всемирного\\[3mm]тяготения}\par
\vspace{8mm}
{\Large Сила тяжести. Ускорение свободного падения}\par
\vspace{23mm}
{\large Персональное пособие для повторения}\par
\vspace{5mm}
{\LARGE\bfseries Володя Хачатурян}\par
\vfill
{\small\color{darkaccent} Теория \textbullet\ Разобранные примеры \textbullet\ Алгоритм \textbullet\ 10 задач}\par
\vspace{12mm}
{\normalsize\bfseries Лёвин Артём Александрович}\\[2mm]
{\normalsize эксклюзивно для Володи Хачатуряна}\par
\vspace{22mm}
\end{center}
\begin{tikzpicture}[remember picture,overlay]
\draw[accent!55,line width=1.6pt]
  ([xshift=17mm,yshift=26mm]current page.south west)
  .. controls ([xshift=72mm,yshift=41mm]current page.south west)
  and ([xshift=141mm,yshift=9mm]current page.south west)
  .. ([xshift=-17mm,yshift=27mm]current page.south east);
\end{tikzpicture}
\end{titlepage}

\sect{1. Закон всемирного тяготения}
\checkline Все тела, имеющие массу, притягиваются друг к другу. Для двух материальных точек (или сферически-симметричных тел вне друг друга) модуль силы:
\[
\boxed{\displaystyle F=G\frac{m_1m_2}{r^2}},
\qquad
G\approx 6{,}67\cdot10^{-11}\ \frac{\text{Н}\cdot\text{м}^2}{\text{кг}^2}.
\]
Здесь $m_1,m_2$ --- массы тел в килограммах, $r>0$ --- \key{расстояние между их центрами} в метрах, $F$ --- сила в ньютонах, $G$ --- гравитационная постоянная.

\begin{center}
\begin{tikzpicture}[scale=0.89,>=Latex]
\tkzDefPoint(0,0){A}
\tkzDefPoint(1,0){C}
\tkzDefPoint(8,0){B}
\tkzDefPoint(8.67,0){D}
\tkzDrawCircle[draw=accent!70,line width=1pt,fill=softblue](A,C)
\tkzDrawCircle[draw=accent!70,line width=1pt,fill=softblue](B,D)
\tkzDrawPoints[size=3,color=darkaccent](A,B)
\draw[-{Latex[length=3mm]},line width=1.6pt,accent] (A)--(1.8,0);
\draw[-{Latex[length=3mm]},line width=1.6pt,accent] (B)--(6.2,0);
\node[above=13mm] at (A) {$m_1$};
\node[above=12mm] at (B) {$m_2$};
\node[above=2mm] at (1.9,0) {$\vec F_{1\leftarrow 2}$};
\node[above=2mm] at (6.1,0) {$\vec F_{2\leftarrow 1}$};
\draw[<->,line width=0.8pt] (0,-1.30)--node[fill=white,inner sep=1.5pt] {$r$} (8,-1.30);
\draw[densely dotted] (A)--(0,-1.45);
\draw[densely dotted] (B)--(8,-1.45);
\end{tikzpicture}
\end{center}
\checkline \key{Третий закон Ньютона:} силы имеют одинаковые модули и противоположные направления, хотя массы тел могут различаться:
\[
\vec F_{1\leftarrow 2}=-\vec F_{2\leftarrow 1}.
\]
\checkline Если одну массу увеличить в $k$ раз, сила возрастёт в $k$ раз. Если \key{расстояние} увеличить в $k$ раз, сила уменьшится в $k^2$ раз.

\subsect{Контроль смысла}
\begin{itemize}
\item Увеличили расстояние вдвое $\Rightarrow$ сила уменьшилась \key{в 4 раза}.
\item Увеличили обе массы вдвое $\Rightarrow$ сила увеличилась \key{в 4 раза}.
\item Силы равны по модулю, однако ускорения тел различаются: $a=F/m$.
\end{itemize}

\sect{2. Сравнение двух сил: главный приём занятия}
\checkline Если у тел один и тот же центральный объект массы $M$, то
\[
F_1=G\frac{Mm_1}{r_1^2},\qquad F_2=G\frac{Mm_2}{r_2^2}.
\]
\checkline Делим \key{в том порядке, который указан в вопросе}, и заменяем деление умножением на обратную дробь:
\[
\frac{F_1}{F_2}
=\frac{GMm_1}{r_1^2}\cdot\frac{r_2^2}{GMm_2}
=\boxed{\frac{m_1}{m_2}\left(\frac{r_2}{r_1}\right)^2}.
\]
Если массы исследуемых тел одинаковы, $m_1=m_2$, остаётся только отношение квадратов расстояний.

\subsect{Пример 1. Два одинаковых спутника}
Радиус первой орбиты втрое больше радиуса второй: $r_1=3r_2$. Найдите $F_2/F_1$.
\[
\frac{F_2}{F_1}=\left(\frac{r_1}{r_2}\right)^2
=\left(\frac{3r_2}{r_2}\right)^2=\boxed{9}.
\]
\key{Вывод:} второй спутник притягивается в 9 раз сильнее.

\subsect{Пример 2. По силам найдите расстояния}
Два тела одинаковой массы притягиваются к одной звезде. Известно: $F_1=4F_2$. Найдите $r_1/r_2$.
\[
\frac{F_1}{F_2}=\left(\frac{r_2}{r_1}\right)^2=4
\quad\Longrightarrow\quad
\frac{r_2}{r_1}=2
\quad\Longrightarrow\quad
\boxed{\frac{r_1}{r_2}=\frac12=0{,}5}.
\]
Расстояния положительны, поэтому при извлечении корня выбираем положительное значение.

\subsect{Пример 3. Земля и Марс относительно Солнца}
Пусть $m_{\text{М}}=0{,}1m_{\text{З}}$ и $r_{\text{СМ}}=1{,}5r_{\text{СЗ}}$. Во сколько раз сила притяжения Земли к Солнцу больше, чем Марса к Солнцу?
\[
\frac{F_{\text{СЗ}}}{F_{\text{СМ}}}
=\frac{m_{\text{З}}}{m_{\text{М}}}
\left(\frac{r_{\text{СМ}}}{r_{\text{СЗ}}}\right)^2
=\frac{m_{\text{З}}}{0{,}1m_{\text{З}}}
\left(\frac{1{,}5r_{\text{СЗ}}}{r_{\text{СЗ}}}\right)^2
=\frac{2{,}25}{0{,}1}=\boxed{22{,}5}.
\]

\subsect{Математические ловушки}
\begin{itemize}
\item «В три раза больше» $\Rightarrow$ $r_2=3r_1$; \key{умножаем}, а не прибавляем $3$.
\item «Половина массы» $\Rightarrow$ $m_2=0{,}5m_1$; сохраняйте одни и те же обозначения во всём решении.
\item При подстановке $r_2=1{,}5r_1$ записывайте \key{$(1{,}5r_1)^2=2{,}25r_1^2$}.
\item $\dfrac ab:\dfrac cd=\dfrac ab\cdot\dfrac dc$, если знаменатели и делитель ненулевые.
\item Проверяйте, \key{какое отношение} спрашивают: $r_1/r_2$ или $r_2/r_1$.
\end{itemize}

\clearpage
\sect{3. Сила тяжести и ускорение свободного падения}
\checkline Сила, с которой планета притягивает тело, называется силой тяжести:
\[
\boxed{F_{\text{тяж}}=mg},
\]
где $m$ --- масса тела, $g$ --- ускорение свободного падения в данном месте. У поверхности Земли для приближённых расчётов часто берут $g\approx 10\ \text{м/с}^2$.

\checkline Направление силы тяжести --- \key{к центру планеты}. Вблизи поверхности Земли это направление называют вертикально вниз. Начальная скорость и угол броска \key{не входят} в формулу $F_{\text{тяж}}=mg$.

\subsect{Как получить формулу для $g$}
Если планета имеет массу $M$ и радиус $R$, а тело находится на высоте $h\ge0$, расстояние от тела до центра равно \key{$r=R+h$}. Приравниваем два выражения одной силы:
\[
mg=G\frac{Mm}{(R+h)^2}
\quad\Longrightarrow\quad
\boxed{g(h)=G\frac{M}{(R+h)^2}}.
\]
Для $h=0$ получаем $g_0=GM/R^2$, а значит
\[
\boxed{\frac{g(h)}{g_0}=\left(\frac{R}{R+h}\right)^2}.
\]
\key{Важно:} высота $h$ отсчитывается от поверхности, расстояние $r$ --- от центра. Модель предполагает сферически-симметричную планету и не учитывает её вращение.

\subsect{Пример 4. Тело брошено под углом}
Камень массой $200$ г брошен под углом $45^\circ$ с начальной скоростью $15$ м/с. Найдите силу тяжести сразу после броска ($g=10$ м/с$^2$).
\[
200\ \text{г}=0{,}2\ \text{кг},\qquad
F_{\text{тяж}}=0{,}2\cdot10=\boxed{2\ \text{Н}}.
\]
Угол и скорость здесь лишние данные. Ответ относится к моменту вблизи поверхности.

\subsect{Пример 5. Космонавт на орбите Луны}
На поверхности Луны сила тяжести равна $180$ Н. Найдите её на расстоянии $1{,}5R$ от \key{центра} Луны (это высота $0{,}5R$ над поверхностью), где $R$ --- радиус Луны.
\[
r_1=R,\quad r_2=1{,}5R,\qquad
F_2=F_1\left(\frac{r_1}{r_2}\right)^2
=\frac{180}{(1{,}5)^2}=\boxed{80\ \text{Н}}.
\]

\clearpage
\begin{center}
{\Large\bfseries\color{darkaccent} Алгоритм: сравнение сил тяготения}\par
\vspace{5mm}
\begin{tikzpicture}
\node[flowterminal] (start) at (0,0) {Задача на две силы \(F_1\) и \(F_2\)};
\node[flowstep] (one) at (0,-2.05) {Переведите текст в формулы: \(m_2=km_1\), \(r_2=qr_1\) и т.\,п.};
\node[flowstep] (two) at (0,-4.45) {Запишите закон тяготения отдельно для каждого случая};
\node[flowstep] (three) at (0,-6.85) {Разделите формулы в нужном порядке; при делении дробей переверните вторую};
\node[flowstep] (four) at (0,-9.25) {Сократите общие множители; сохраните квадраты расстояний};
\node[flowdecision] (choice) at (0,-11.94) {Ищете отношение расстояний?};
\node[flowbranch] (left) at (-4.42,-15.25) {\textbf{Да}\par Выразите квадрат отношения \(r_1/r_2\) или \(r_2/r_1\); извлеките положительный корень};
\node[flowbranch] (right) at (4.42,-15.25) {\textbf{Нет}\par Подставьте данные в скобках; возведите множители в квадрат и вычислите};
\node[flowterminal] (end) at (0,-18.86) {Сверьте порядок отношения, единицы и физический смысл ответа};
\draw[flowarrow] (start)--(one);
\draw[flowarrow] (one)--(two);
\draw[flowarrow] (two)--(three);
\draw[flowarrow] (three)--(four);
\draw[flowarrow] (four)--(choice);
\draw[flowarrow] (choice.west) -| (left.north);
\draw[flowarrow] (choice.east) -| (right.north);
\node[fill=white,inner sep=1pt,font=\small] at (-2.60,-12.60) {да};
\node[fill=white,inner sep=1pt,font=\small] at (2.60,-12.60) {нет};
\draw[draw=darkaccent,line width=0.9pt] (left.south) -- (-4.42,-17.58) -- (0,-17.58);
\draw[draw=darkaccent,line width=0.9pt] (right.south) -- (4.42,-17.58) -- (0,-17.58);
\draw[flowarrow] (0,-17.58) -- (end.north);
\end{tikzpicture}
\end{center}

\clearpage
\sect{4. Самостоятельная тренировка \textnormal{(Путь А)}}
\textbf{10 заданий от базового уровня к повышенному.} В заданиях 1--2 используйте $g=10$ м/с$^2$. Во всех остальных считайте тела сферически-симметричными, а расстояния --- измеренными от центров, если явно не сказано иное.
\vspace{2mm}
\begin{enumerate}[label=\textbf{\arabic*.},itemsep=10pt]
\item Найдите силу тяжести, действующую у поверхности Земли на тело массой $3$ кг.
\item Камень массой $250$ г брошен со скоростью $12$ м/с под углом $60^\circ$ к горизонту. Каков модуль силы тяжести сразу после броска?
\item Два тела притягиваются друг к другу с силой $F$. Как изменится сила, если массу \emph{каждого} тела увеличить вдвое, сохранив расстояние между центрами?
\item Два одинаковых спутника обращаются вокруг одной планеты; $r_1=4r_2$. Найдите отношение $F_2/F_1$.
\item Для двух одинаковых спутников одной планеты известно: $F_1=25F_2$. Найдите $r_1/r_2$.
\item На поверхности Луны на тело действует сила тяжести $270$ Н. Какая сила действует на него на расстоянии $3R$ от центра Луны? Здесь $R$ --- радиус Луны.
\item Два тела обращаются вокруг одной планеты. Масса второго составляет половину массы первого: $m_2=0{,}5m_1$. При этом $r_1=2R$, $r_2=8R$. Найдите $F_1/F_2$.
\item Две планеты притягиваются к одной звезде. Масса второй составляет четверть массы первой, а расстояние второй до звезды вдвое больше расстояния первой. Во сколько раз сила притяжения первой планеты к звезде больше силы притяжения второй?
\item У поверхности планеты $g_0=9{,}6$ м/с$^2$. Найдите ускорение свободного падения на высоте $h=R$ над её поверхностью, где $R$ --- радиус планеты.
\item На какой высоте $h$ над поверхностью планеты ускорение свободного падения станет в $16$ раз меньше, чем на поверхности? Ответ выразите через радиус планеты $R$.
\end{enumerate}

\vfill
\textcolor{darkaccent}{\textbf{Перед сдачей проверьте:}} единицы СИ, скобки при подстановке, квадраты расстояний и порядок отношения.

\clearpage
\sect{5. Ответы к тренировочному блоку}
\renewcommand{\arraystretch}{1.42}
\begin{center}
\begin{tabularx}{0.93\linewidth}{@{}c X@{}}
\toprule
\textbf{№} & \textbf{Ответ} \\
\midrule
1 & $30$ Н \\
2 & $2{,}5$ Н \\
3 & Увеличится в $4$ раза: $F_{\text{нов}}=4F$ \\
4 & $F_2/F_1=16$ \\
5 & $r_1/r_2=1/5=0{,}2$ \\
6 & $30$ Н \\
7 & $F_1/F_2=32$ \\
8 & В $16$ раз \\
9 & $2{,}4$ м/с$^2$ \\
10 & $h=3R$ \\
\bottomrule
\end{tabularx}
\end{center}
\vspace{9mm}
\subsect{Короткая самопроверка}
\begin{itemize}
\item В заданиях 4--8 используйте одну и ту же структуру:
\[
\frac{F_1}{F_2}=\frac{m_1}{m_2}\left(\frac{r_2}{r_1}\right)^2
\quad\text{(общая центральная масса).}
\]
\item В заданиях 9--10 применяйте $g(h)/g_0=(R/(R+h))^2$.
\item Ускорение и сила тяжести уменьшаются с удалением от центра планеты; проверьте это по знаку результата.
\end{itemize}
\vfill
\begin{center}\small\textcolor{darkaccent}{Лёвин Артём Александрович эксклюзивно для Володи Хачатуряна}\end{center}
\end{document}
